\section{字符串}
\subsection{字符串哈希}
$\mathcal O(|S|)$
\begin{lstlisting}[caption={}]
typedef std::pair<i64,i64> PII;
const int N = 1e6 + 10;

i64 pow1[N],pow2[N];
struct string_hash{
    const int Base = 499,M1 = 998244853,M2 = 989244353; 
    std::string str;
    std::vector<i64> hash1,hash2;
    void extend(char c,int idx){
        pow1[idx] = (pow1[idx - 1] * Base) % M1;
        hash1[idx] = (hash1[idx - 1] * Base + c) % M1;
        pow2[idx] = (pow2[idx - 1] * Base) % M2;
        hash2[idx] = (hash2[idx - 1] * Base + c) % M2;
    }

    void init(int n){
        hash1.resize(n + 1,0);
        hash2.resize(n + 1,0);
		for(int i = 0;i <= n + 2;i++)
			pow1[i] = pow2[i] = 1;
    }

    void init(std::string str){
        i64 n = str.size();
        init(n);
        this->str = str;
        for(int i = 0;i < n;i++){
            extend(str[i],i + 1);
        }
    }
    string_hash(){}
    string_hash(std::string str){init(str);}

    PII getHash(i64 l,i64 r){ // 1_idx
        i64 t1 = (hash1[r] - 1LL * hash1[l - 1] * pow1[r - l + 1]) % M1;
        t1 = (t1 + M1) % M1;
        i64 t2 = (hash2[r] - 1LL * hash2[l - 1] * pow2[r - l + 1]) % M2;
        t2 = (t2 + M2) % M2;
        return {t1,t2};
    }

    bool same(i64 l1,i64 r1,i64 l2,i64 r2){
        PII ta = getHash(l1,r1);
        PII tb = getHash(l2,r2);
        return ta.first == tb.first && ta.second == tb.second;
    }
};
\end{lstlisting}

\subsection{最小表示法}
最小表示法是一种用于解决字符串最小表示问题的方法。它的核心思想是找到一个字符串的所有循环同构字符串中字典序最小的那个字符串,$\mathcal O(|S|)$。
\begin{lstlisting}[caption={}]
template<typename T>
int get_min_str(const T& str){  //返回循环同构字符串的最小字典序首位下标
    int i = 0,j = 1,k = 0;
    int n = str.size();
    while(i < n && j < n && k < n){
        int tmp = str[(i + k) % n] - str[(j + k) % n];
        if(tmp == 0){
            k++;
        }
        else{
            if(tmp > 0) i += k + 1;
            else j += k + 1;
            if(i == j) j++;
            k = 0;
        }
    }
    return std::min(i,j);
}
\end{lstlisting}


\subsection{KMP}
在字符串中查找子串,以最坏 $\mathcal O(N+M)$ 的时间计算 $t$ 在 $s$ 中出现的全部位置。

最小周期：字符串长度-整个字符串的 $border$;

最小循环节：区别于周期，当字符串长度 $n \bmod (n - nxt[n]) = 0$ 时，等于最小周期，否则为 $n$ 。
\begin{lstlisting}[caption={}]
i64 ne[N],m,n;
std::string a,b;
std::vector<i64> ans;
void getnext(std::string& str){
	ne[0] = 0;
	for(int i = 1,j = 0;i < str.size();i++){
		while(j && str[i] != str[j]) j = ne[j - 1];
		if(str[i] == str[j]) j++;
		ne[i] = j;
	}
}

void kmp(std::string& a,std::string& b){
	getnext(b);
	for(int i = 0,j = 0;i < a.size();i++){
		while(j && a[i] != b[j]) j = ne[j - 1];
		if(a[i] == b[j]) j++;
		if(j == b.size()) ans.push_back(i - b.size() + 2);
	}
}
\end{lstlisting}

\subsection{Z函数(扩展KMP)}
获取字符串 $s$ 和 $s[i,n-1]$ (即以 $s[i]$ 开头的后缀)
的最长公共前缀(LCP)的长度,总复杂度 $\mathcal O(N)$。
\begin{lstlisting}[caption={}]
template<typename T>
std::vector<int> z_function(const T& str){
    int n = str.size();
    std::vector<int> z(n);
    int j = 0;
    for(int i = 1;i < n;i++){
        z[i] = (j + z[j] <= i) ? 0 : std::min(j + z[j] - i,z[i - j]);
        while(i + z[i] < n && str[z[i]] == str[i + z[i]]){
            z[i]++;
        }
        if(j + z[j] < i + z[i]){
            j = i;
        }
    }
    z[0] = n;
    return z;
}
\end{lstlisting}

\subsection{Manacher}
$\mathcal O(|S|)$
\begin{lstlisting}[caption={}]
struct Manacher{
    std::vector<int> p;
    std::string change(std::string& str){
        std::string temp = "!#";
        for(int i = 0;i < str.size();i++){
            temp += str[i];
            temp += '#';
        }
        return temp;
    }

    void get_p(std::string& str){
        int maxright = 0,mid = 0;
        str = change(str);
        p.assign(str.size(),0);
        for(int i = 1;i < str.size();i++){
            if(i < maxright){
                p[i] = std::min(p[2 * mid - i],maxright - i + 1);
            }
            else{
                p[i] = 1 + (str[i] != '#');
            }
            while(i - p[i] >= 0 && p[i] + i < str.size() && str[i - p[i]] == str[i + p[i]]){
                p[i]++;
            }
            if(i + p[i] - 1 > maxright){
                mid = i;
                maxright = i + p[i] - 1;
            }
        }
    }
    // 返回最大回文串长度
    int get_max(std::string& str){
        get_p(str);
        int res = 0;
        for(int i = 1;i < p.size();i++){
            res = std::max(res,p[i] - 1);
        }
        return res;
    }
};
\end{lstlisting}


\subsection{字典树}
$\mathcal O(N \log N)$
\begin{lstlisting}[caption={}]
template<int N = 26>
struct Trie{
    std::vector<std::array<int,N>> son;
    std::vector<int> cnt;

    Trie(){
        son.push_back(std::array<int,N>{});
        cnt.push_back(0);
    }

    void insert(std::string& str){  //插入 
        int p = 0,n = str.size();
        for(int i = 0;i < n;i++){
            int ch = str[i] - '0';
            if(!son[p][ch]){ //如果此处字母不存在就创建一个
                son[p][ch] = son.size();  
                son.push_back(std::array<int,N>{});
                cnt.push_back(0);
            }
            p = son[p][ch];  //指向下一个位置 
            // cnt[p]++;
        }
        cnt[p]++;  //该处单词数+1 
    }

    int find(std::string& str){  //找 
        int p = 0,n = str.size();
        for(int i = 0;i < n;i++){
            int ch = str[i] - '0';
            if(!son[p][ch]){
                return 0;
            }
            p = son[p][ch];
        }
        return cnt[p];
    }
};
\end{lstlisting}

\subsection{AC自动机}
$\mathcal O(\sum|s_i|+|S|)$
\begin{lstlisting}[caption={}]
struct ACAM{
    vector<int> ne,in,cnt,f;
    vector<vector<int>> son;
    int idx = 0;
    ACAM(){
        son.resize(N,vector<int>(26,0));
        ne.resize(N,0),in.resize(N,0);
        cnt.resize(N,0),f.resize(N,0);
    }
    ll insert(string str){  //插入模式串
        int p = 0;
        for(int i = 0;str[i];i++){
            int a = str[i] - 'a';
            if(!son[p][a]) son[p][a] = ++idx;  //如果此处字母不存在就创建一个 
            p = son[p][a];  //指向下一个位置 
        }
        cnt[p]++;  //该处单词数+1 
        return p;
    }

    void build(){
        queue<int> q;
        for(int i = 0;i < 26;i++)
            if(son[0][i]) q.push(son[0][i]);
        while(!q.empty()){
            int u = q.front();
            q.pop();
            for(int i = 0;i < 26;i++){
                int v = son[u][i];
                if(v){
                    ne[v] = son[ne[u]][i];  //建立回跳边
                    q.push(v);
                    in[son[ne[u]][i]]++;
                }
                else son[u][i] = son[ne[u]][i];  //建转移边
            }
        }
    }
    void query(string str,int n){
        for(int k = 0,i = 0;k < str.size();k++){
            i = son[i][str[k] - 'a'];
            f[i]++;
            // for(int j = i;j && ~cnt[j];j = ne[j]){
            //     ans += cnt[j];
            //     cnt[j] = -1;//不重复计数
            // }
        }
        queue<ll> q;
        for(int i = 1;i <= idx;i++)
            if(in[i] == 0) q.push(i);
        while(!q.empty()){
            ll u = q.front(),v = ne[u];
            q.pop();
            f[v] += f[u],in[v]--;
            if(v != 0 && in[v] == 0) q.push(v);
            f[u] *= cnt[u];
        }
    }
};
\end{lstlisting}

\subsection{后缀数组}
把字符串的所有后缀按字典序排序后得到的$rk$数组,以 $\mathcal O(N \log{N})$ 的复杂度求解。
\begin{lstlisting}[caption={}]
struct SA{
	std::string str;
	int n,m,tp = 128,p = 0;
	std::vector<int> rk,sa;
	SA(std::string _s){
		str = " " + _s;
		n = _s.size();
		m = (n << 1) + 1;
		rk.assign(m,0),sa.assign(m,0);
		work();
	};
	void work(){
		if(n == 1){
			sa[1] = rk[1] = 1;
			return;
        }
		std::vector<int> cnt(std::max(tp,n) + 1),id(n + 1);
		for(int i = 1;i <= n;i++) cnt[rk[i] = str[i]]++;									
		for(int i = 1;i <= tp;i++) cnt[i] += cnt[i - 1];
		for(int i = n;i >= 1;i--) sa[cnt[rk[i]]--] = i;

		for(int w = 1;w < n;w <<= 1,tp = p){
			int cur = 0;
			for(int i = n - w + 1;i <= n;i++) id[++cur] = i;
			for(int i = 1;i <= n;i++)
				if(sa[i] > w) id[++cur] = sa[i] - w;
			std::fill(cnt.begin(),cnt.end(),0);
			for(int i = 1;i <= n;i++) cnt[rk[i]]++;
			for(int i = 1;i <= tp;i++) cnt[i] += cnt[i - 1];
			for(int i = n;i >= 1;i--) sa[cnt[rk[id[i]]]--] = id[i];
			std::vector<int> temp(m);
			p = 0;
			for(int i = 1;i <= n;i++){
				if(rk[sa[i]] != rk[sa[i - 1]] || rk[sa[i] + w] != rk[sa[i - 1] + w]) p++;
				temp[sa[i]] = p; 
			}
			rk = temp;
			if(p == n) break;
		}
	}
};
\end{lstlisting}

\subsection{回文自动机 PAM(回文树)}
本质不同的回文串个数：$idx - 2$;回文子串出现次数。

对于一个字符串 $s$，它的本质不同回文子串个数最多只有 $|s|$ 个,
那么,构造 $s$ 的回文树的时间复杂度是 $\mathcal O(|s|)$ 。
\begin{lstlisting}[caption={}]
struct PalindromeAutomaton {
    constexpr static int N = 5e5 + 10;
    int tr[N][26], fail[N], len[N];
    int cntNodes, last;
    int cnt[N];
    string s;
    PalindromeAutomaton(string s) {
        memset(tr, 0, sizeof tr);
        memset(fail, 0, sizeof fail);
        len[0] = 0, fail[0] = 1;
        len[1] = -1, fail[1] = 0;
        cntNodes = 1;
        last = 0;
        this->s = s;
    }
    void insert(char c, int i) {
        int u = get_fail(last, i);
        if (!tr[u][c - 'a']) {
            int v = ++cntNodes;
            fail[v] = tr[get_fail(fail[u], i)][c - 'a'];
            tr[u][c - 'a'] = v;
            len[v] = len[u] + 2;
            cnt[v] = cnt[fail[v]] + 1;
        }
        last = tr[u][c - 'a'];
    }
    int get_fail(int u, int i) {
        while (i - len[u] - 1 <= -1 || s[i - len[u] - 1] != s[i]) {
            u = fail[u];
        }
        return u;
    }
};
\end{lstlisting}